%%%%%%%%%%%%%%%%%%%%%%%%%%%%%%%%%%%%%%%
% Cover Letter - BETA Technologies
% Software, Controls & Autonomy Internship (Summer 2027)
% Plain XeLaTeX — matches simple resume style (Spectral font)
%%%%%%%%%%%%%%%%%%%%%%%%%%%%%%%%%%%%%%%

\documentclass[11pt]{article}

\usepackage[margin=0.75in, top=0.7in, bottom=0.7in]{geometry}
\usepackage[hidelinks]{hyperref}
\usepackage{parskip}
\usepackage{fontspec}

\setmainfont{Spectral}[
  Path = ../../cv/fonts/,
  UprightFont    = Spectral-Regular.ttf,
  BoldFont       = Spectral-Bold.ttf,
  ItalicFont     = Spectral-Italic.ttf,
  BoldItalicFont = Spectral-BoldItalic.ttf
]

\setlength{\parindent}{0pt}
\setlength{\parskip}{14pt}
\linespread{1.15}
\raggedright

% No hyphenation
\hyphenpenalty=10000
\exhyphenpenalty=10000

\pagestyle{empty}
\urlstyle{same}

\begin{document}

% ── DATE ─────────────────────────────────────────────────────────────────────
\begin{flushright}
  \today
\end{flushright}

\vspace{4pt}

% ── BODY ─────────────────────────────────────────────────────────────────────
Dear BETA Technologies Hiring Team,

I am applying for the Software, Controls \& Autonomy Internship for Summer
2027, specifically the Embedded Systems and Airborne Electric Software
tracks. I'm a Computer Science student at RIT in the combined bachelor's and
master's program, with an expected graduation date of December 2027, and
have a strong interest in firmware and embedded software. At GE Vernova, I
built a full I/O subsystem from scratch in C on an STM32F7 for industrial
radios, using the STM32 Hardware Abstraction Layer. On the RIT Electric
Vehicle Team, I'm the vehicle control unit (VCU) Software Lead. The firmware
I've built so far is for our current track motorcycle's VCU, written in both
bare-metal C++ and ThreadX RTOS, running a 300\,Hz control loop for
drive-by-wire control and board-level fault handling, and integrating a
Simulink-generated state machine from our controls engineers into that
firmware. In my free time I'm building
an autonomous fixed-wing flight controller on an ESP32S3 running FreeRTOS,
with PID stabilization, sensor fusion, GPS waypoint following, and manual
override for safety. The autonomous side is validated in simulation, and the
full airframe recently completed its first flight in manual mode.

The ALIA VTOL's lift-plus-cruise layout is what got me interested in BETA
specifically. Instead of tilting rotors or variable-pitch props to
transition between hover and cruise, the propellers stay fixed and the wing
carries cruise flight, trading a harder hover-to-cruise transition for a
mechanically simpler aircraft. That's exactly the kind of engineering
choice I find interesting, and the fourfold motor redundancy on each lift
prop is the same fault-tolerant thinking I try to build into the VCU. I want
to help write the embedded software behind decisions like that.

The VCU's firmware design is informed by ISO 26262 principles, and I built
its CAN/CANopen communication layer from scratch, including the CANopen
object dictionary and message formats connecting the VCU to its motor
controller, battery management system, and power subsystem, which is close
to the safety-critical electrical software work your Airborne Electric
Software team is describing. I've worked across SPI, I2C, UART, ADC, and PWM building
both the GE Vernova I/O subsystem and EVT's peripheral driver library, and I
debug at the hardware level regularly with a Saleae logic analyzer, an
oscilloscope, and a J-Link Segger. My flight controller project is the
closest thing I have to aircraft work directly. It integrates the Madgwick
AHRS library for IMU and magnetometer sensor fusion and attitude state
estimation, runs a cascaded PID
stabilization and waypoint-following loop, and I designed its PCB in KiCad.

I would welcome the opportunity to talk more about where I could contribute.
Thank you for your consideration.

Thanks,

Travis Brown

\end{document}
